\section{Camera Triggering} \label{sec:camera_triggering}

\subsection{Trajectory generation}
Having finished with the most critical part of the inspection, locating the wind turbine in the local reference frame, it is necessary to plan the trajectory with regards to its position. The generated trajectories allow the UAV drone to capture images of points of interest on the rotating blades.

The generation process involves a Python script which produces trajectories and stores them in a .csv file format. The script is designed to generate a trajectory relative in position to the wind turbine, taking into account the specific parameters of the wind turbine itself. These are its base height and rotor radius. The Python script's flexibility allows for easy adjustment and adaptation through parametrization, ensuring that the drone can navigate the inspection process effectively and efficiently on virtually any model of the turbine.

Through the generated trajectory, the inspection encompasses the front and back sides of the rotor blade seen from three different angles. To capture comprehensive visual information, the drone's trajectory should go as follows:
\begin{enumerate}
    \item Approach the rotor center head-on at a safe distance.
    \item Move sideways to the tip of the blades.
    \item Make a 180\degree turn around the tip to reach the back side of the rotor.
    \item Move sideways to the rotor center behind the turbine.
    \item Repeat with a height offset above and below the traversed line, pointing the gimbal at 45 degrees.
\end{enumerate}

\begin{figure}
    \centering
    \includegraphics[width=0.45\textwidth]{figures/trajectory_model.png}
    \caption{Generated rotating wind turbine model with the generated trajectory}
    \label{fig:trajecotry_model}
\end{figure}

Figure~\ref{fig:trajecotry_model} depicts the generated point cloud of the 3d model representing the wind turbine and the trajectory generated with the adequate parameters.


This trajectory generation approach enhances the efficiency and accuracy of wind turbine inspections by streamlining the data collection process. The collected images can subsequently be analyzed and processed for analysis, such as detecting potential defects or monitoring the condition of the wind turbine.



\subsection{Camera Triggering}

To capture images of the rotating blades during autonomous wind turbine inspection using a UAV drone, a precise camera triggering process has to be created.

The image capture process has to be synchronized with the passing of the blades. To ensure that images are taken at the right time, the camera triggering mechanism contains a camera and a laser range finder fitted onto a gimbal, an example of which is given in~\ref{fig:gimbal}. This gimbal is used to point the two in the desired direction and an Arduino micro-controller is used to connect the laser range finder to the camera trigger.


As the blade approaches while the UAV is waiting in position, the laser range finder detects the blade's presence in front of it, signaling the Arduino micro-controller to activate the camera trigger. Upon receiving the signal from the laser range finder, the camera trigger initiates the image capture process, capturing a picture of the rotor blade from the desired angle. This synchronized approach guarantees that the images correspond to the targeted points of interest, providing valuable data for further analysis.


\begin{figure}
    \centering
    \includegraphics[width=0.45\textwidth]{figures/inspection_simulation/Schematics-of-a-3-axes-gimbal-system-Courtesy-of-DJI.png}
    \caption{An example of a gimbal system for a commercial UAV drone.  Source image from \cite{Fu-Gimbal}.}
    \label{fig:gimbal}
\end{figure}